% !PRINTED 427
\weberpage{427}
Um die \gesperrt{Galois}'sche Gruppe dieser Gleichung $6^{\mathrm{ten}}$ Grades zu erhalten, haben wir den Einfluss zu untersuchen, den die Anwendung der Ikosa\"edersubstitutionen auf das System der Nebengruppen (2) hat. Es gen\"ugt dazu, die Substitutionen $\Theta$, $\psi$, $\chi$ zu betrachten. Es ist aber, wenn wir die Reihenfolge der Nebengruppen in (2) beachten, nach \S.\,58, (23), (25):
\[
\begin{aligned}
P\Theta &= Q + Q\chi\Theta + Q\chi\Theta^2 + Q\chi\Theta^3 + Q\chi\Theta^4 + Q\chi \\
P\psi &= Q + Q\chi + Q\chi\Theta^4 + Q\chi\Theta^3 + Q\chi\Theta^2 + Q\chi\Theta \\
P\chi &= Q\chi + Q + Q\chi\Theta + Q\chi\Theta^3 + Q\chi\Theta^2 + Q\chi\Theta^4,
\end{aligned}
\]
und daraus folgt, dass sich die Indices von $U$ folgendermaassen vertauschen:
\[
\begin{array}{c@{\quad}cccccc}
 & \infty, & 0, & 1, & 2, & 3, & 4 \\
\cline{2-7}
\Theta) & \infty & 1 & 2 & 3 & 4 & 0 \\
\psi) & \infty & 0 & 4 & 3 & 2 & 1 \\
\chi) & 0 & \infty & 1 & 3 & 2 & 4.
\end{array}
\]
Bezeichnen wir den Index allgemein mit $\xi$, und nehmen, wie im siebenten Abschnitte, $\xi$ nach dem Modul 5, so dass zwei nach dem Modul 5 congruente Zahlen als nicht verschieden gelten, so k\"onnen wir diese Vertauschungen so darstellen:
\[
\Theta = (\xi,\ \xi+1),\quad \psi = (\xi,\ -\xi),\quad \chi = \left(\xi,\ \frac{1}{\xi}\right)\!.
\]
\gesperrt{Daraus folgt aber, dass die Gruppe unserer Gleichung} $6^{\mathrm{ten}}$ \gesperrt{Grades mit der Congruenzgruppe} $L_5$ (\S.\,68) \gesperrt{\"ubereinstimmt.}

Um zur Bildung von Resolventen $6^{\mathrm{ten}}$ Grades zu gelangen, m\"ussen wir \"ahnlich wie bei den Resolventen $5^{\mathrm{ten}}$ Grades verfahren.

Wir haben schon im \S.\,55 die Grundformen der Di\"edergruppen kennen gelernt. Wir stellen sie jetzt in der Form dar:
\[
\varphi_1 = x_1x_2,\quad \varphi_2 = x_1^5+ix_2^5,\quad \varphi_3 = x_1^5-ix_2^5.
\]
Nehmen wir $\Theta$ und $\psi$ mit der Determinante 1, also
\[
\Theta = \begin{pmatrix} e^{\frac{\pi i}{5}}, & 0 \\ 0, & e^{-\frac{\pi i}{5}} \end{pmatrix},\quad \psi = \begin{pmatrix} 0, & 1 \\ -1, & 0 \end{pmatrix};
\]
so \"andern sich die Grundformen $\varphi_1, \varphi_2, \varphi_3$ in folgender Weise:
\[
\begin{array}{c@{\quad}ccc}
 & \varphi_1 & \varphi_2 & \varphi_3 \\
\cline{2-4}
\Theta) & \varphi_1 & -\varphi_2 & -\varphi_3 \\
\psi) & -\varphi_1 & -i\varphi_2 & i\varphi_3.
\end{array}
\]
